\chapter{Sostenibilità Sociale}

\section{Installazione di csdetector}
L'installazione e la configurazione di \texttt{csdetector} richiede alcuni passaggi per preparare l'ambiente e configurare le dipendenze necessarie. Di seguito è riportata una guida passo passo su come è stato installato:

\begin{enumerate}
    \item \textbf{Clonazione della repository} di \texttt{csdetector} da GitHub:
    \begin{lstlisting}[language=bash, breaklines=true]
git clone https://github.com/Nuri22/csdetector.git
cd csdetector
    \end{lstlisting}
    
    \item \textbf{Installazione delle dipendenze Python} necessarie utilizzando \texttt{pip}:
    \begin{lstlisting}[language=bash, breaklines=true]
pip install -r requirements.txt
    \end{lstlisting}
    
    \item \textbf{Configurazione dell'ambiente}: Lo strumento richiede un file di configurazione per memorizzare le credenziali. Copia il file \texttt{.env.example} fornito per creare un file \texttt{.env}:
    \begin{lstlisting}[language=bash, breaklines=true]
cp .env.example .env
    \end{lstlisting}
\end{enumerate}

\section{Limitazioni dell'Analisi}
Per eseguire un'analisi completa del repository utilizzando \texttt{csdetector}, è strettamente necessario configurare il file \texttt{.env} con un Fine-Grained Personal Access Token di GitHub. Questo token consente allo strumento di recuperare i commit, le issue e le pull request per costruire un grafo delle interazioni tra i contributori e identificare successivamente potenziali community smells.

Tuttavia, siccome questa l'applicazione web 2chance era originariamente il mio progetto di Software Dependability, ma il proprietario del repository è un'altra persona con cui ho collaborato, per motivi di sicurezza, non ho ricevuto l'autorizzazione per generare o utilizzare un fine-grained token con i permessi necessari sulla sua repository. Di conseguenza, non ho potuto analizzare la cronologia dei commit precedenti per identificare se fossero presenti dei community smells durante la collaborazione originale.

Inoltre, per quanto riguarda la fase attuale del progetto, non è stato possibile effettuare alcuna analisi perché ci ho lavorato da solo. Poiché i community smells derivano intrinsecamente dalle dinamiche del team, dai modelli di comunicazione e dalle interazioni tra più sviluppatori, l'esecuzione di un'analisi dei community smells su un progetto singolo non produrrebbe risultati. Pertanto, l'analisi con \texttt{csdetector} non è stata effettuata per questo progetto.
